% Part: methods
% Chapter: proofs
% Section: resources

\documentclass[../../../include/open-logic-section]{subfiles}

\begin{document}

\olfileid{mth}{prf}{res}

\section{इतर साधने}

गणितात सिद्धता कशा कराव्यात याविषयी अनेक उपयुक्त पुस्तके आहेत. मूळ स्रोताने विशेषतः \emph{How to Read and do Proofs: An Introduction to Mathematical Thought Processes} \citep{Solow2013} आणि \emph{How to Prove It: A Structured Approach} \citep{Velleman2019} पाहण्याची सूचना केली आहे. \href{http://www.people.vcu.edu/~rhammack/BookOfProof/BookOfProof.pdf}{\emph{Book of Proof}} \citep{Hammack2013} आणि \href{https://scholarworks.gvsu.edu/books/7/}{\emph{Mathematical Reasoning}} \citep{Sandstrum2019} ही सिद्धतांवरील पुस्तके मूळ स्रोतामध्ये इंटरनेटवर मुक्तपणे उपलब्ध म्हणून नमूद केलेली आहेत. तत्त्वज्ञानाच्या विद्यार्थ्यांसाठी \emph{More Precisely: The Math you need to do Philosophy} \citep{Steinhart2018} हे गणितीय विचाराचे उपयुक्त प्रास्ताविक पुस्तक ठरू शकते.

इंटरनेटवर सिद्धतांविषयी काही छोटी मार्गदर्शकेही आहेत. उदाहरणार्थ, \href{https://math.berkeley.edu/~hutching/teach/proofs.pdf}{“Introduction to Mathematical Arguments”} \citep{Hutchings2003} आणि \href{https://eugeniacheng.com/wp-content/uploads/2017/02/cheng-proofguide.pdf}{“How to write proofs”} \citep{Cheng2004}.

\subsection{प्रेरणादायी चित्रफिती}

गृहपाठ करण्याची प्रेरणाच उरली नाही असे वाटते का? मन खिन्न आहे का? या चित्रफिती मदत करू शकतील!

\begin{itemize}
\item \url{https://www.youtube.com/watch?v=ZXsQAXx_ao0}
\item \url{https://www.youtube.com/watch?v=BQ4yd2W50No}
\item \url{https://www.youtube.com/watch?v=StTqXEQ2l-Y}
\end{itemize}

\end{document}
